\appendix
\section{Upper-Level Substitution Bias}
\label{app:substitution-bias}

This appendix describes the upper-level substitution-bias diagnostic used to
compare the package's baseline group-specific HICP aggregation with a chained
Tornqvist alternative. The purpose of the exercise is not to replace the
official HICP convention, but to quantify how sensitive measured inflation
inequality is to the upper-level aggregation formula used to combine COICOP
price relatives.

\subsection{Conceptual background}

In a fixed-basket price index, expenditure weights are held fixed while prices
change. Such an index does not fully account for substitution across
consumption categories when relative prices change. If households shift
expenditure away from categories whose prices rise relatively faster, a
fixed-weight index will generally grow faster than a symmetric index that uses
information on both initial and later expenditure shares. This difference is
usually interpreted as an upper-level substitution bias.

The Bureau of Labor Statistics applies this logic by comparing a research CPI
based on a fixed-basket formula with a chained superlative index. In the HICP
setting considered here, the reference index is kept closer to the Eurostat
convention: monthly COICOP price relatives are aggregated with annual
expenditure weights and then chain-linked over time. The comparison index is a
chained Tornqvist index computed from the same matched HICP-HBS data.

\subsection{Index definitions}

Let \(p_{j,t}\) denote the elementary price index for COICOP category \(j\) in
month \(t\). Let \(r_{j,t}\) denote the corresponding unchained price relative
used in the monthly aggregation. For household group \(g\), let
\(w_{j,g,y}\) be the annual expenditure weight assigned to COICOP category
\(j\) in year \(y\). These weights are obtained by combining official HICP
item weights with HBS expenditure shares by income group.

For each household group \(g\), the package first computes a Laspeyres-type
monthly aggregate using the annual weights available for the corresponding
link period:

\[
L_{g,t}
= \sum_j s_{j,g,y(t)} r_{j,t},
\]

where \(s_{j,g,y(t)}\) denotes the normalized expenditure share for group
\(g\), category \(j\), and the annual weight year associated with month \(t\).
The monthly aggregates are then chain-linked to obtain a Eurostat-style
chained Laspeyres index:

\[
P^{CL}_{g,t}
= \operatorname{chain}\left(L_{g,t}\right).
\]

The comparison index uses the Tornqvist formula at the same upper level. Let
\(s_{j,g,y(t)-1}\) and \(s_{j,g,y(t)}\) denote the previous-year and current-year
shares used for the link. The monthly Tornqvist aggregate is:

\[
T_{g,t}
= \prod_j r_{j,t}^{\frac{1}{2}
\left(s_{j,g,y(t)-1} + s_{j,g,y(t)}\right)}.
\]

These Tornqvist aggregates are also chain-linked:

\[
P^{CT}_{g,t}
= \operatorname{chain}\left(T_{g,t}\right).
\]

\subsection{Annualized growth rates}

For each group \(g\), the comparison is made using annualized growth rates over
the selected sample window. If \(t_0\) and \(t_1\) are the first and last months
available in the comparison, and \(m=t_1-t_0\) is the number of elapsed months,
the annualized growth rate of index \(P\) is:

\[
\gamma_g(P)
= \left[
\left(\frac{P_{g,t_1}}{P_{g,t_0}}\right)^{12/m} - 1
\right] \times 100.
\]

The upper-level substitution-bias diagnostic is then defined as:

\[
\operatorname{SB}_g
= \gamma_g\left(P^{CL}\right)
- \gamma_g\left(P^{CT}\right).
\]

In the package output, these terms are reported as:

\[
\texttt{substitution\_bias}
= \texttt{chained\_laspeyres\_growth}
- \texttt{chained\_toernqvist\_growth}.
\]

A positive value indicates that the chained Laspeyres aggregation grows faster
than the chained Tornqvist alternative over the same period.

\subsection{Example for France}

Table~\ref{tab:fr-substitution-bias} reports the diagnostic for France by
income quintile, using Eurostat HICP and HBS data at COICOP level 2 from 2019
onwards. The calculation corresponds to:

\begin{verbatim}
calculate_substitution_bias(
  "FR",
  "income",
  level = 2,
  start_year = 2019
)
\end{verbatim}

\begin{table}[htbp]
\centering
\caption{Upper-level substitution bias by income group, France, COICOP level 2}
\label{tab:fr-substitution-bias}
\begin{tabular}{lrrr}
\hline
Income group & Chained Laspeyres growth & Chained Tornqvist growth & Substitution bias \\
\hline
Total & 2.96 & 2.90 & 0.06 \\
First quintile & 2.83 & 2.74 & 0.10 \\
Second quintile & 2.94 & 2.85 & 0.10 \\
Third quintile & 2.98 & 2.90 & 0.09 \\
Fourth quintile & 3.01 & 2.93 & 0.08 \\
Fifth quintile & 2.99 & 2.92 & 0.07 \\
\hline
\end{tabular}
\begin{minipage}{0.9\textwidth}
\footnotesize
\emph{Notes:} Growth rates and substitution-bias estimates are expressed in
annualized percentage points over the selected sample window. The chained
Laspeyres index follows the package's HICP-style aggregation with annual
weights. The chained Tornqvist index uses previous-year and current-year
category weights in the upper-level aggregation.
\end{minipage}
\end{table}

\subsection{Interpretation and limitations}

The estimates in Table~\ref{tab:fr-substitution-bias} are small but positive
for all income groups. This indicates that, over the selected period, the
Eurostat-style chained Laspeyres aggregation produces slightly higher
annualized price growth than the chained Tornqvist alternative. The difference
is somewhat larger for the first income quintile than for the fifth quintile in
this example, suggesting that the choice of upper-level formula can affect the
measured distributional profile of inflation, even when the magnitude remains
modest.

Several limitations should be kept in mind. First, Eurostat HBS expenditure
shares by income group are not observed monthly. The Tornqvist comparison is
therefore approximated with annual category-specific weights rather than with
high-frequency household expenditure data. Second, HBS waves are available only
periodically and are matched to annual HICP weights using the package's HBS-CPI
timing rule. Third, the calculation is performed at the selected COICOP level;
results may differ when more detailed national HBS sources are used. The
diagnostic should therefore be read as a robustness check on the aggregation
formula, not as an official chained consumer price index.
